% Part: proof-theory
% Chapter: proof-search
% Section: tableaux

\documentclass[../../../include/open-logic-section]{subfiles}

\begin{document}

\olfileid{pt}{ps}{tab}

\olsection{تابلوګانې}

د تابلو سيستمونه د !!{proof} هغه سيستمونه دي چې د
سيکوېنټ حسابونو سره نژدې تړاو لري او د ثبوت د لټون دپاره،
هم په لاس او هم په سافټوير کښې، ځانګړې آسانتيا لري۔ په
تابلو کښې د سيکوېنټونو د ونې پر ځاے !!a{proof} د
!!{formula}s ونه ده۔ خو د طبيعي استنتاج پر خلاف ئې قواعد د
سيکوېنټ حساب د قواعدو په شان دي۔ د تابلو !!{proof} په اصل
کښې د يوه مدل څرګند لټون دے؛ ځکه کله کله ورته
\emph{معنايي} تابلو يا \emph{د صدق ونه} هم وايي۔

\emph{نښه لرونکے} تابلو د نښه لرونکو !!{formula}s ونه
ده، چې $\sFmla{\True}{!A}$ يا $\sFmla{\False}{!A}$ بڼه
لرلے شي۔ ونې ښکته لور ته غځېږي۔ په هغو !!{formula}s پيلېږي
چې زموږ د !!{prove} کولو مقصد ټاکي۔ مثلاً، د سيکوېنټ
$!A, !B \Sequent !C, !D$ د !!{prove} کولو دپاره به په
$\sFmla{\True}{!A}, \sFmla{\True}{!B},
\sFmla{\False}{!C}, \sFmla{\False}{!D}$ پيل کېږي۔ داسې
!!^a{structure} چې $!A$ او $!B$ رښتيا او $!C$ او $!D$ دروغ
کړي، هغه !!a{structure} دے چې ښيي
$!A, !B \Sequent !C, !D$ معتبر نۀ دے۔
\textbf{د سرچينه‌يي نښې سمون OLCMP-195:} د ښي اړخ وروستي
فارمول ناسمه رښتيا نښه په دروغ نښه بدله شوې، څو د ټول ښي
اړخ د ناسمولو هماغه شرط پوره شي۔

د تابلو پاڼه، يعنې تر ټولو لاندې غوټه، د څانګې د
غځولو په قواعدو غځېدے شي: په هماغه څانګه نورې غوټې زياتوي
يا لاندې دوې هم‌مورې غوټې ورزياتوي او څانګه په دوو نويو
څانګو وېشي۔ قواعد په هماغه څانګه کښې پر هر پاسني نښه
لرونکي !!{formula} کارېدلے شي۔ د کلاسيکي تابلو سيستم
\Log{Tc} د څانګې د غځولو قواعد په \olref{tab:Tc} کښې دي۔

\subfile{rules-Tc}

لکه چې ښکاري، قواعد هماغه د \Log{G3c} قواعد دي خو سرچپه
شوي؛ $\True$ او $\False$ نښې په ترتيب ښيي چې
!!a{formula} د سيکوېنټ په کيڼ يا ښي اړخ اصلي فارمول دے۔

د تابلو يوه څانګه هله \emph{تړلې} ده چې دواړه
$\sFmla{\True}{!A}$ او $\sFmla{\False}{!A}$ ولري؛ يا
$\sFmla{\True}{\lfalse}$ ولري۔ که هره څانګه تړلې وي، ټول
تابلو تړلے بولو۔
\textbf{د سرچينه‌يي تړلو د شرط سمون OLCMP-198:} په رښتيا
نښه لرونکے دروغ هم څانګه تړي؛ د سيکوېنټ د کيڼ دروغ
بديهيت همدا حالت دے، چې په اصل کښې نۀ و ذکر شوے۔ د لاندې
تابلو بېلګه تړلې ده۔
مثلاً، د $!A \lor (!B \land !C) \Sequent (!A \lor !B) \land
(!A \lor !C)$ دپاره دلته يو تړلے تابلو دے؛ يعنې په
$\sFmla{\True}{!A \lor (!B \land
!C)}$ او $\sFmla{\False}{(!A \lor !B) \land (!A \lor !C)}$
پيلېږي:
\begin{oltableau}
  [\sFmla{\True}{\formula{A} \lor (\formula{B} \land \formula{C})},
    just=\TAss, checked
    [\sFmla{\False}{(\formula{A} \lor \formula{B}) \land
        (\formula{A} \lor \formula{C})}, just=\TAss, checked
      [\sFmla{\True}{\formula{A}}, just={\TRule{\True}{\lor}[1]}
        [\sFmla{\False}{\formula{A} \lor \formula{B}},
          just={\TRule{\False}{\land}[2]}, checked
          [\sFmla{\False}{\formula{A}},
            just={\TRule{\False}{\lor}[4]}
            [\sFmla{\False}{\formula{B}},
              just={\TRule{\False}{\lor}[4]}, close]
          ]
        ]
        [\sFmla{\False}{\formula{A} \lor \formula{C}},
          just={\TRule{\False}{\land}[2]}, checked
          [\sFmla{\False}{\formula{A}},
            just={\TRule{\False}{\lor}[4]}
            [\sFmla{\False}{\formula{C}},
              just={\TRule{\False}{\lor}[4]}, close]
          ]
        ]
      ]
      [\sFmla{\True}{\formula{B} \land \formula{C}},
        just={\TRule{\True}{\lor}[1]}, checked
        [\sFmla{\False}{\formula{A} \lor \formula{B}},
          just={\TRule{\False}{\land}[2]}, checked
          [\sFmla{\True}{\formula{B}},
            just={\TRule{\True}{\land}[3]}, move by=2
            [\sFmla{\True}{\formula{C}},
              just={\TRule{\True}{\land}[3]}
              [\sFmla{\False}{\formula{A}},
                just={\TRule{\False}{\lor}[4]}
                [\sFmla{\False}{\formula{B}},
                  just={\TRule{\False}{\lor}[4]},close
                ]
              ]
            ]
          ]
        ]
        [\sFmla{\False}{\formula{A} \lor \formula{C}},
          just={\TRule{\False}{\land}[2]}, checked
          [\sFmla{\True}{\formula{B}},
            just={\TRule{\True}{\land}[3]}, move by=2
              [\sFmla{\True}{\formula{C}},
                just={\TRule{\True}{\land}[3]}
                [\sFmla{\False}{\formula{A}},
                  just={\TRule{\False}{\lor}[4]}
                  [\sFmla{\False}{\formula{C}},
                  just={\TRule{\False}{\lor}[4]},close
                ]
              ]
            ]
          ]
        ]
      ]
    ]
  ]
\end{oltableau}

د کتنې نښې ښيي چې اړونده قاعده په !!a{formula} باندې
په ټولو لاندې څانګو کښې کارول شوې ده۔ د $\otimes$ نښه د
څانګې تړل ښيي۔ د قواعدو ورپسې د کرښو شمېرې ښيي چې قاعده
په کومو نښه لرونکو !!{formula}s کارول شوې ده؛ يعنې په
هماغه څانګه کښې پر هغه !!{formula} چې ښودل شوې کرښه ئې لري۔

په تابلو کښې د ثبوت لټون آسان دے۔ تابلو په پړاوونو
جوړوو۔ په پړاو $0$ کښې پيليز !!{formula}s پاس وليکئ۔ په
پړاو $n+1$ کښې د هرې نۀ تړلې پاڼې غوټې دپاره داسې وکړئ:
لومړے، يعنې تر ټولو پاسنے، ناعټومي نښه لرونکے فارمول ومومئ
چې لا د کتنې نښه نۀ لري او د لاندې تکراري کميت ټاکنې له
نوعې نۀ وي۔ په دغه فارمول د څانګې د غځولو قاعده وکاروئ۔
کله چې قاعده په هره څانګه کښې وکارول شي چې همدغه نښه
لرونکے فارمول لري، د کتنې نښه پرې کېږدئ۔ (پورته بېلګه په
همدې الګوريتم جوړه شوې ده۔)
\textbf{د راکمولو د نوعې سمون OLCMP-196:} اټومي فارمول
د دغو غځولو قواعدو هېڅ قاعده نۀ لري؛ يوازې ناعټومي مورد
ټاکل کېږي، او تړلې څانګه نوره نۀ غځېږي۔

د $\sFmla{\True}{\lforall}$ او
$\sFmla{\False}{\lexists}$ بڼې نښه لرونکي !!{formula}s
هېڅکله د کتنې نښه نۀ اخلي، نو ځانګړې کارونه غواړي۔ که
موجود وي، بايد يقيني کړو چې $\TRule{\True}{\forall}$ او
$\TRule{\False}{\lexists}$ قواعد بالاخره پر ټولو د کارونې
وړ ترمونو وکارول شي۔ مثلاً، دا په دې تضمينېدے شي چې هر
پړاو کښې په هره څانګه پر دغو ټولو فارمولونو قاعده وکاروو،
وروسته له دې چې د دغې بڼې پرته تر ټولو پاسنے ناعټومي نښه
لرونکے !!{formula} وکارول شي۔ د کارونې ترم $t$ هغه لومړے
ترم~$t$ دے چې په څانګه کښې له مخکې موجودو !!{constant}s
(يا $\Obj
c_0$، که هېڅ !!{constant} موجود نۀ وي) او په پيليزو
!!{formula}s کښې له راغلو !!{function}s څخه جوړېدے شي، او
اړوند نښه لرونکے !!{formula} $\sFmla{\True}{!B(t)}$ يا
~$\sFmla{\False}{!B(t)}$ لا په څانګه کښې موجود نۀ وي۔
\textbf{د ترمونو د اشباع سمون OLCMP-197:} که په دغه
پړاو کښې داسې نااستعمال شوے ترم نۀ وي، همدغه قاعده نوې
غوټه نۀ زياتوي؛ خو د کتنې نښه لا نۀ اخلي، څو د څانګې له
وروستيو نويو ثابتونو سره بيا کارېدے شي۔

که د ثبوت دا لټون تړلے تابلو جوړ نۀ کړي، نو يا يوه
متناهي پرانيستې مشبوع څانګه لري چې ټول عادي ناعټومي
!!{formula}s ئې د کتنې نښې لري او تکراري کميت ټاکنه ئې پر
ټولو اوسني د کارونې وړ ترمونو کارول شوې ده؛ يا لټون يوه
نامتناهي، متناهي څانګېدونکې ونه جوړوي چې بايد يوه نامتناهي
څانګه ولري۔ لکه په \olref[cpl]{sec} کښې، داسې
!!a{structure}~$\Struct{M}$ تعريفولے شو چې که
$\sFmla{\True}{!A}$ په څانګه کښې راځي، $\Sat{M}{!A}$، او
که $\sFmla{\False}{!A}$ راځي، $\Sat/{M}{!A}$ وي۔
($\Theta =
\Setabs{!A}{\sFmla{\True}{!A} \text{ په څانګه کښې}}$ او
$\Xi =
\Setabs{!A}{\sFmla{\False}{!A} \text{ په څانګه کښې}}$
واخلئ۔)
\textbf{د بشپړتيا د حد څرګندونه OLCMP-197/200:} متناهي
څانګه د تکراري کميت ټاکنې له نښې پرته هم مشبوع کېدے شي۔
د OLCMP-190 په ترمي مدل کښې د ازادو متغيرونو د ګومارنې
جلا پاتې دليل دلته هم پاتې دے؛ عمومي معنايي ادعا ساتل
شوې، خو د ازادو متغيرونو د ټولو حالتونو نوے ثبوت نۀ دے
ورزيات شوے۔

د ثبوت د لټون په دې طريقه جوړ شوي تابلوګانې هرې غوټې
ته د سيکوېنټ په ورکولو په آسانۍ د \Log{G3c} په ثبوتونو
بدلېدے شي۔ که پيليز فارمولونه د سيکوېنټ
$\Gamma \Sequent
\Delta$ سره برابر وي، هر پيليز فارمول په
$\Gamma \Sequent
\Delta$ ونوموئ۔ که څانګه پر نښه لرونکي
!!{formula} $\sFmla{\True}{!A}$ د غځولو د قاعدې په کارونه
وغځېږي، د پاڼې نوم $!A, \Pi \Sequent \Lambda$ دے؛ او که
پر نښه لرونکي !!{formula} $\sFmla{\False}{!A}$ کارول شوې
وي، د پاڼې نوم $\Pi \Sequent
\Lambda, !A$ دے۔ هر چې تابلو د \TRule{\True}{\star} قاعده
کاروي، پر سيکوېنټ د \LeftR{\star} قاعده وکاروئ، چې $!A$
ئې اصلي فارمول وي؛ او هر چې د \TRule{\False}{\star} قاعده
کاروي، د \RightR{\star} قاعده وکاروئ۔ د قاعدې مقدمې هغه
سيکوېنټونه دي چې په تابلو کښې ورزياتې شوې اړوندې غوټې
نوموي۔ کله چې تابلو په $\sFmla{\True}{!A}$ او
$\sFmla{\False}{!A}$ وتړل شي، د وروستۍ پاڼې نوم به
$!A, \Pi \Sequent \Lambda,
!A$ بڼه ولري۔ د تابلو ځينې پرله‌پسې غوټې به يو شان
سيکوېنټ واخلي؛ تکراري موردونه لرې کړئ۔
\textbf{د رسمي بديهيت د حد OLCMP-199، د OLCMP-191 تعقيب:}
د ګډ اټومي فارمول تړل په G3c کښې مستقيم بديهيت ورکوي، او
په رښتيا نښه لرونکے دروغ د کيڼ دروغ بديهيت ورکوي۔ د
مرکب ګډ فارمول تړل معنايي صحيح دي، خو د هويت جلا مشتق
ثبوت غواړي، يوازې بې‌مقدمې رسمي G3c پاڼه نۀ ده۔ لاندې
رسم د خپلو درې برخو د اټومي لوستلو په صورت کښې رسمي
G3c ثبوت دے؛ د عمومي مرکب هويت نوې ثبوتونه نۀ دي ورزيات
شوي۔ مثلاً، پورته تابلو په دغه !!{proof} بدلېږي:
\[\small
\Axiom$!A \fCenter !A, !B$
\RightLabel{\RightR{\lor}}
\UnaryInf$!A \fCenter !A \lor !B$
\Axiom$!A \fCenter !A, !C$
\RightLabel{\RightR{\lor}}
\UnaryInf$!A \fCenter !A \lor !C$
\RightLabel{\RightR{\land}}
\BinaryInf$!A \fCenter (!A \lor !B) \land (!A \lor !C)$
\Axiom$!B, !C \fCenter !A, !B$
\RightLabel{\RightR{\lor}}
\UnaryInf$!B, !C \fCenter !A \lor !B$
\RightLabel{\LeftR{\land}}
\UnaryInf$!B \land !C \fCenter !A \lor !B$
\Axiom$!B, !C \fCenter !A, !C$
\RightLabel{\RightR{\lor}}
\UnaryInf$!B, !C \fCenter !A \lor !C$
\RightLabel{\LeftR{\land}}
\UnaryInf$!B \land !C \fCenter !A \lor !C$
\RightLabel{\RightR{\land}}
\BinaryInf$!B \land !C \fCenter (!A \lor !B) \land
(!A \lor !C)$
\RightLabel{\LeftR{\lor}}
\BinaryInf$!A \lor (!B \land !C) \fCenter (!A \lor !B) \land
(!A \lor !C)$
\DisplayProof
\]

\end{document}
